\documentclass[10pt]{article}

% ---------- Document Identity (update these only) ----------
\newcommand{\DocStage}{}
\newcommand{\DocTitle}{Your Road to Tier One SOF}
\newcommand{\ProgramName}{182-Week Civilian-to-Selection Readiness Pipeline}
\newcommand{\ProgramSubtitle}{}

% ---------- Page ----------
\usepackage[legalpaper,left=0.5in,right=0.5in,top=0.6in,bottom=0.45in]{geometry}

% ---------- Typography ----------
\usepackage[T1]{fontenc}
\usepackage[utf8]{inputenc}
\usepackage{libertinus}
\usepackage{microtype}
\usepackage{parskip}
\usepackage{ragged2e}
\usepackage{textcomp}
\AtBeginDocument{\color{black}}
\AtBeginDocument{\pagecolor{white}}
\AtBeginDocument{\justifying}

% ---------- Landscape pages ----------
\usepackage{pdflscape}

% ---------- Color ----------
\usepackage[table]{xcolor}
\definecolor{StageBlue}{HTML}{1F4E79}
\definecolor{StageBlueDark}{HTML}{163A5A}
\definecolor{LightGray}{HTML}{F2F4F7}
\definecolor{MidGray}{HTML}{D9DEE5}
\definecolor{RuleGray}{HTML}{C7CED8}
\definecolor{HeaderInk}{HTML}{000000}
\definecolor{Ink}{HTML}{000000}

% ---------- Utility ----------
\usepackage{enumitem}
\sloppy
\setlength{\emergencystretch}{2em}

% ---------- Boxes / UI ----------
\usepackage[most]{tcolorbox}

\tcbset{
  charterTitle/.style={
    enhanced,
    colback=StageBlue!4,
    colframe=StageBlue,
    boxrule=1.25pt,
    arc=3pt,
    left=16pt,right=16pt,top=14pt,bottom=14pt,
    borderline north={3.2pt}{0pt}{StageBlue},
    borderline west={4pt}{0pt}{StageBlue},
    borderline south={1.25pt}{0pt}{StageBlue},
    drop shadow={black!25!white},
  },
  charterRule/.style={
    enhanced,
    colback=LightGray,
    colframe=RuleGray,
    boxrule=0.85pt,
    arc=3pt,
    left=12pt,right=12pt,top=10pt,bottom=10pt,
    borderline west={3pt}{0pt}{StageBlue},
  },
  charterBadge/.style={
    enhanced,
    colback=StageBlue,
    coltext=white,
    colframe=StageBlueDark,
    boxrule=0.6pt,
    arc=2pt,
    left=6pt,right=6pt,top=3pt,bottom=3pt,
  }
}
\newtcbox{\Badge}{nobeforeafter,charterBadge}

% ---------- Lists ----------
\setlist[itemize]{leftmargin=1.5em, itemsep=0.25em, topsep=0.25em}
\setlist[enumerate]{leftmargin=1.8em, itemsep=0.25em, topsep=0.25em}

% ---------- TOC ----------
\usepackage{tocloft}
\setcounter{tocdepth}{3}
\setlength{\cftbeforesecskip}{3pt}
\setlength{\cftbeforesubsecskip}{2pt}
\setlength{\cftbeforesubsubsecskip}{0pt}
\renewcommand{\cftsecfont}{\bfseries}
\renewcommand{\cftsecpagefont}{\bfseries}
\renewcommand{\cftsubsecfont}{\normalfont}
\renewcommand{\cftsubsecpagefont}{\normalfont}
\renewcommand{\cftsubsubsecfont}{\normalfont}
\renewcommand{\cftsubsubsecpagefont}{\normalfont}
\setlength{\cftsecindent}{0pt}
\setlength{\cftsubsecindent}{1.25em}
\setlength{\cftsubsubsecindent}{2.8em}
\setlength{\cftsecnumwidth}{2.1em}
\setlength{\cftsubsecnumwidth}{2.8em}
\setlength{\cftsubsubsecnumwidth}{3.6em}

% Remove the built-in TOC heading so only the custom heading is shown
\makeatletter
\renewcommand{\tableofcontents}{\@starttoc{toc}}
\makeatother

% ---------- Header/Footer: page number only ----------
\usepackage{fancyhdr}
\pagestyle{fancy}
\fancyhf{}
\renewcommand{\headrulewidth}{0pt}
\renewcommand{\footrulewidth}{0pt}
\fancyfoot[C]{\thepage}

% ---------- Section formatting ----------
\usepackage{titlesec}

\titleformat{\section}
  {\Large\bfseries\color{Ink}}
  {\thesection}{0.6em}{}
  [\vspace{0.25em}\textcolor{StageBlue}{\rule{\linewidth}{0.8pt}}]

\titleformat{\subsection}
  {\large\bfseries\color{Ink}}
  {\thesubsection}{0.6em}{}

\titleformat{\subsubsection}
  {\normalsize\bfseries\color{Ink}}
  {\thesubsubsection}{0.6em}{}

\titlespacing*{\section}{0pt}{0.95em}{0.35em}
\titlespacing*{\subsection}{0pt}{0.65em}{0.25em}
\titlespacing*{\subsubsection}{0pt}{0.55em}{0.2em}

% ---------- Table engine ----------
\usepackage{tabularray}
\UseTblrLibrary{booktabs}

% ---------- tabularray: GLOBAL table treatment ----------
\SetTblrInner{
  cells = {fg=black, bg=white, valign=t},
}

% ---------- Header helpers ----------
\newcommand{\Hdr}[1]{\shortstack[c]{\strut #1\strut}}

% ---------- Lines ----------
\newcommand{\TitleLine}{\textcolor{StageBlue}{\rule{0.98\linewidth}{0.95pt}}}
\newcommand{\SubTitleLine}{\textcolor{RuleGray}{\rule{0.98\linewidth}{0.65pt}}}

% ---------- Continued on next page ----------
\newcommand{\ContinuedOnNextPageInline}{%
  \par
  \vfill
  \noindent\textcolor{RuleGray}{\rule{\linewidth}{1pt}}\vspace{-2pt}
  \noindent\hfill\textit{Continued on next page}\par\vspace{-10pt}
  \noindent\textcolor{RuleGray}{\rule{\linewidth}{1pt}}
  \clearpage
}

% ---------- PDF hyperlinks ----------
\usepackage[hidelinks]{hyperref}
\hypersetup{
  colorlinks=false,
  pdfborder={0 0 0},
  linktoc=all,
  pdftitle={\DocTitle},
  pdfauthor={\ProgramName}
}

% =========================================================
\begin{document}

\begin{tcolorbox}
\begin{center}
Stage 3 — Test Calendar, Deload Map, and Macro Stress Pattern\\[0.35em]
SOF Physical Development Transformation Program\\[0.25em]
182-Week Civilian-to-Selection Readiness Pipeline\\[0.65em]
\TitleLine
\end{center}
\end{tcolorbox}

\vspace{0.35em}

\section*{Stage 3 — Test Calendar, Deload Map, and Macro Stress Pattern}
\phantomsection
\addcontentsline{toc}{section}{Stage 3 — Test Calendar, Deload Map, and Macro Stress Pattern}

\subsection*{Introduction}
\phantomsection
\addcontentsline{toc}{subsection}{Introduction}

Stage 3 exists to schedule testing and deload windows so gate evidence is produced safely, under comparable conditions, and with admissibility protected across the full 182-week pipeline (Phase 00 through Phase 13).

Stage 3 controls when tests may be executed as gate-grade events and when fatigue-reduction periods occur so:

\begin{itemize}
\item gate-critical tests are not executed under accumulated fatigue that contaminates comparability,
\item high-consequence domains (\textbf{RUN}, \textbf{RUCK}, \textbf{SWIM}, underwater governance) are not forced into unsafe conditions by calendar pressure,
\item replication requirements (2 of last 3 \textbf{VALID} for gate-tier classification) are achievable without “overtesting” or chronic fatigue accumulation,
\item invalid tests are handled by reslotting, not by compression or stacking,
\item environmental constraints in Grand Forks, North Dakota (ice, wind, cold, heat, low daylight, variable access) are treated as predictable schedule constraints rather than exceptions.
\end{itemize}

Division of authority is non-negotiable:

\begin{itemize}
\item Stage 2 controls how tests are executed and what makes them \textbf{VALID} or \textbf{INVALID} (protocol requirements, environment/tool logging, evidence standards, and replication rules).
\item Stage 3 controls when tests are executed as gate-grade and how deload/test windows are spaced and reslotted.
\item Stage 3 \textbf{MUST} NOT prescribe training programming, workouts, sets/reps, intervals, pacing plans, or day-by-day session construction.
\end{itemize}

Core enforcement outcome:

\begin{itemize}
\item A test executed outside an authorized Stage 3 protected window may be logged informationally, but it \textbf{MUST} be treated conservatively for gate use when fatigue, novelty, environment risk, or incomplete logging is present. Gate-grade use is restricted to protected windows by default doctrine.
\end{itemize}

\subsubsection*{2) Scope}

\paragraph*{2.1 Inclusions}

Stage 3 includes enforceable scheduling doctrine that preserves safety and validity.

\begin{enumerate}
\item Macro stress doctrine (work-to-deload rhythm)
\begin{itemize}
\item Defines a repeating work-to-deload rhythm used across phases to prevent chronic fatigue accumulation and to protect test validity.
\item Establishes the default posture: deload week equals test week.
\item Defines the “protected window” concept: test-grade evidence is intended to be produced inside deload/test windows.
\item Defines minimum time separation between protected windows to prevent continuous testing posture.
\end{itemize}

\item Protected deload/test windows doctrine (gate-grade test authorization)
\begin{itemize}
\item Defines the only authorized windows for gate-grade testing:
  \begin{itemize}
  \item Mid-phase checkpoint window (when phases are long enough to warrant a mid-phase checkpoint).
  \item End-phase gate window (required for gate outcomes when tests are specified as gate inputs).
  \end{itemize}
\item Defines what “protected” means operationally:
  \begin{itemize}
  \item Reduced fatigue posture is required.
  \item No novelty posture applies (environment/tool/footwear changes are restricted).
  \item Anti-stacking spacing rules apply (see section 2.1.4).
  \end{itemize}
\item Defines eligibility for attempting gate-grade tests:
  \begin{itemize}
  \item prerequisites and readiness are governed by Stage 1.F admissibility and Stage 2.C safety screens,
  \item durability and stop-rule constraints override schedule.
  \end{itemize}
\end{itemize}

\item Default doctrine: deload week equals test week
\begin{itemize}
\item The deload/test window is a combined:
  \begin{itemize}
  \item fatigue reduction period, and
  \item validity protection period.
  \end{itemize}
\item Within that week, high-cost tests are scheduled by spacing doctrine, not by convenience.
\item Testing does not “replace” training. Stage 0 cadence remains 6 sessions/week; the deload/test week preserves cadence while reducing fatigue and protecting comparability (Stage 3 controls timing; programming methods remain out-of-scope).
\end{itemize}

\item Anti-stacking and spacing posture (doctrine-level only)

Stage 3 defines spacing rules to avoid test contamination and injury risk. These rules apply within each protected window and when reslotting.

Test event classes (for spacing doctrine; no new IDs; references are domain categories only):

\begin{itemize}
\item High-cost locomotion: \textbf{RUN} time trials (\textbf{C8}–\textbf{C11}), \textbf{RUCK} time trials (\textbf{C12}), \textbf{SWIM} time trials (\textbf{C13}), underwater governance-constrained exposure (\textbf{C14} when authorized).
\item High-cost calisthenics: \textbf{CAL} max/timed events (\textbf{C4}–\textbf{C7}).
\item High-cost strength proxy: \textbf{STR} proxy tests (\textbf{C17}) when executed as formal test events.
\item High-cost non-run conditioning: \textbf{AER} time trials (\textbf{C18}) when executed as formal test events.
\item Governance/trend markers: \textbf{BODYCOMP} measures (\textbf{C1}–\textbf{C3}), \textbf{RECOVERY} markers (\textbf{C15}), \textbf{DURABILITY} screen flags (\textbf{C16}).
\end{itemize}

Spacing doctrine (minimum enforceable rules):

\begin{itemize}
\item No more than one high-cost locomotion test event per 48 hours inside a protected window.
\item No more than one high-cost “total-system” day (locomotion test day plus additional major test) inside a protected window.
\item \textbf{CAL} events may be grouped only when:
  \begin{itemize}
  \item artifact requirements can be satisfied without compromising auditability,
  \item the grouping does not create a “fatigue cascade” into locomotion tests within the same window.
  \end{itemize}
\item Underwater (\textbf{C14}) is scheduled only if governance prerequisites are already satisfied; underwater is never used to “fill a slot” in a test week.
\item If spacing cannot be satisfied within the protected window, the correct posture is partial completion and reslotting remaining tests to the next protected window; stacking is prohibited.
\end{itemize}

\item Environment reslot/substitute posture for Grand Forks, North Dakota

Stage 3 treats seasonal hazards and access variability as expected constraints:

\begin{itemize}
\item Unsafe outdoor conditions trigger indoor substitution only when Stage 2 validity fields preserve comparability (e.g., treadmill admissibility with full settings logged).
\item If an indoor substitute cannot preserve the construct’s validity (e.g., \textbf{SWIM} without a verified pool; underwater without supervision; unverified distance), the test is reslotted.
\item Test scheduling must assume:
  \begin{itemize}
  \item winter ice and low daylight risk for outdoor routes,
  \item heat and air quality risk in warmer months,
  \item wind and visibility volatility,
  \item facility access volatility.
  \end{itemize}
\end{itemize}

\item Adjustment posture (missed tests, invalid tests, illness, regression, risk flags)

Stage 3 defines schedule-level responses:

\begin{itemize}
\item Missed/invalid tests are reslotted, not compressed.
\item Illness and stop-rule events trigger conservative reslotting and may force a “no testing until stable” posture.
\item Regression and Medical Hold posture override test schedules; calendar does not force testing.
\end{itemize}

Detailed adjustment rules are owned by Stage 3.E, but Stage 3 establishes the controlling principle: reslot rather than compress; preserve comparability; preserve safety.

\item Replication enablement posture (2 of last 3 \textbf{VALID} without overtesting)

Stage 3 must provide sufficient cadence opportunities for Gate-Critical metrics to meet replication requirements without encouraging excessive testing.

Replication enablement rules (schedule posture):

\begin{itemize}
\item Gate-critical metrics must have recurring protected windows such that:
  \begin{itemize}
  \item repeated \textbf{VALID} tests can be achieved across consecutive windows,
  \item failure in one window does not force unsafe “make-up testing.”
  \end{itemize}
\item Replication is not “three tests in one week.” Replication is achieved across windows under comparable conditions.
\item When a decision window is referenced for replication, Stage 3 defines that “metric evidence window” may span multiple protected windows as needed to locate the last three \textbf{VALID} events while still respecting Stage 1.F gate posture and Stage 3 reslot rules.
\end{itemize}
\end{enumerate}

\paragraph*{2.2 Exclusions}

Stage 3 explicitly excludes:

\begin{itemize}
\item Workouts, sets/reps, intervals, weekly splits, microcycles, daily session design, pace charts, and specific training prescriptions.
\item New tests, thresholds, domain tags, metric IDs, protocol IDs, validity tokens, or schema changes. Stage 2 is authoritative.
\item Equipment acquisition guidance and procurement planning. Stage 6–Stage 7 own equipment timing and build-out.
\item Medical diagnosis or treatment. Only stop posture and escalation references are permitted.
\end{itemize}

\subsubsection*{3) How to Use}

Stage 3 is used by Coach Lead and Trainee as the timing control layer that prevents invalid or unsafe evidence from driving gates.

\paragraph*{3.1 Authorized testing posture}

\begin{itemize}
\item Gate-grade tests are executed only inside protected deload/test windows.
\item Tests executed outside protected windows may be logged informationally, but are presumed non-gate-grade unless:
  \begin{itemize}
  \item the event meets all Stage 2 validity fields,
  \item fatigue confounders are controlled,
  \item novelty is absent,
  \item and a Coach Lead explicitly records that the event is treated as gate-grade (rare; requires strong justification; still must obey spacing doctrine).
  \end{itemize}
\end{itemize}

\paragraph*{3.2 Protected window doctrine: what "deload/test week" means operationally}
Protected window is not a license to test everything. It is a controlled evidence production period.

Minimum protected window characteristics:

\begin{itemize}
\item Reduced fatigue posture relative to prior work weeks is required so tests represent capacity rather than exhaustion.
\item High-cost test density is capped by spacing doctrine.
\item Novelty is minimized:
  \begin{itemize}
  \item avoid new route/machine/pool,
  \item avoid new footwear or ruck configuration,
  \item avoid new timing method,
  \item avoid new environment unit IDs immediately before gate-grade testing unless forced (if forced, treat as comparability break and default reslot for gate-grade use).
  \end{itemize}
\end{itemize}

\paragraph*{3.3 Reslot rather than compress}

\begin{itemize}
\item If tests are missed or \textbf{INVALID}, the schedule response \textbf{MUST} be reslot to the next authorized protected window.
\item No “make-up testing” is permitted that stacks multiple high-cost tests into a narrow timeframe or forces testing under compromised conditions.
\item If only a subset of the planned test battery can be executed validly within a window:
  \begin{itemize}
  \item log the executed tests,
  \item explicitly mark unexecuted tests as \textbf{NOT EXECUTED} (do not leave silent gaps),
  \item reslot remaining tests to the next protected window.
  \end{itemize}
\end{itemize}

\paragraph*{3.4 Environment substitution posture (indoor vs reslot)}

Stage 3 authorizes indoor substitution only when Stage 2 validity and comparability can be preserved.

Examples of authorized substitution posture (construct-preserving):

\begin{itemize}
\item \textbf{RUN} time trials may use treadmill alternatives when permitted by Stage 2 protocol conditions and full device settings are logged (\textbf{TM}-2E-###; grade recorded; adopted grade setting recorded; speed profile recorded; pauses handled per protocol).
\item \textbf{RUCK} time trials may use treadmill alternatives when permitted by Stage 2 protocol conditions, including dry load verification, \textbf{TM}-2E-###, and full settings logging.
\end{itemize}

Examples of reslot-required posture (construct not preservable):

\begin{itemize}
\item \textbf{SWIM} tests require verified pool length and unit; if pool is unavailable or unverifiable, reslot rather than substitute with non-aquatic evidence.
\item Underwater (\textbf{C14}) is prohibited without certified lifeguard plus dedicated safety spotter; absence forces reslot and “not authorized” logging posture.
\item Any run/ruck test where distance cannot be verified (\textbf{MEASURE} failure) is reslot-required for gate-grade use.
\end{itemize}

\paragraph*{3.5 Replication enablement without overtesting}

\begin{itemize}
\item Replication for gate-tier classification is achieved across protected windows, not by stacking repeated attempts in a single week.
\item If a metric’s replication status is 1/3 at the end of a protected window:
  \begin{itemize}
  \item do not force immediate retest outside the window,
  \item reslot the next attempt to the next protected window.
  \end{itemize}
\item When a phase gate relies on tier classification:
  \begin{itemize}
  \item stage the gate timing so that at least two protected windows have occurred since the last major method/environment change, enabling replicated \textbf{VALID} results under comparable conditions.
  \end{itemize}
\end{itemize}

\subsubsection*{4) Inputs and Assumptions}

\paragraph*{4.1 Inputs}

\begin{itemize}
\item Stage 0 constraints and environment posture, including Grand Forks, North Dakota seasonal hazards and indoor substitution necessity.
\item Stage 1.F governance, including:
  \begin{itemize}
  \item Session Validity Tag: \textbf{VALID} / \textbf{MODIFIED} / \textbf{INVALID}
  \item Test Validity Tag: \textbf{VALID} / \textbf{INVALID}
  \item Gate outcomes: \textbf{ADVANCE} / \textbf{HOLD} / \textbf{REGRESS} / \textbf{MEDICAL} \textbf{HOLD}
  \item Stop posture and escalation dominance
  \item Freeze windows and change control discipline
  \end{itemize}
\item Stage 2 measurement layer:
  \begin{itemize}
  \item Domain tags: \textbf{RUN}, \textbf{SWIM}, \textbf{CAL}, \textbf{RUCK}, \textbf{STR}, \textbf{BODYCOMP}, \textbf{DURABILITY}, \textbf{RECOVERY}, \textbf{AER}
  \item Metric IDs: \textbf{MET-2B}-###
  \item Protocol cards: \textbf{C1}–\textbf{C18} (with applicable readiness-tier and validity rules)
  \item Tier thresholds: \textbf{ENTRY} / \textbf{INTERMEDIATE} / \textbf{SELECTION-READY}
  \item Stage 2.E logging conventions, Environment Unit IDs, and tool-identity posture
  \end{itemize}
\item Phase structure: Phase 00–Phase 13 and fixed durations per upstream macro design.
\end{itemize}

\paragraph*{4.2 Assumption Ledger (only if needed)}

\begin{itemize}
\item At least one safe outdoor route OR a stable indoor device exists to execute \textbf{RUN} and \textbf{RUCK} test constructs when a protected window is reached; when neither exists safely, default posture is reslot.
\item Pool access is not assumed globally and is treated as phase-gated; \textbf{SWIM} tests occur only when pool length/units can be verified and logged.
\item A timing method exists for test windows; if timing capture is unavailable or not logged, the test result is \textbf{INVALID} for gate use.
\end{itemize}

\subsubsection*{5) Interfaces}

\paragraph*{5.1 Upstream dependencies}

Stage 3 depends on and must remain consistent with:

\begin{itemize}
\item Stage 0 identity, operating constraints, safety/legal/medical boundaries, and environment posture.
\item Stage 1.F governance: validity tags, stop posture, freeze windows, and change control.
\item Stage 2 measurement architecture: domain taxonomy, metric inventory, protocol cards, tier thresholds, and logging conventions.
\end{itemize}

\paragraph*{5.2 Downstream consumers}

Stage 3 serves as the schedule foundation for:

\begin{itemize}
\item Stage 4–Stage 5 phase specifications, which consume protected windows to determine mid-phase checkpoints and end-phase gate timing.
\item Stage 8 crosswalk audits, which check Stage 3 schedules for compatibility with:
  \begin{itemize}
  \item Stage 2 test instruments,
  \item Stage 6–Stage 7 resource feasibility,
  \item environment identity consistency.
  \end{itemize}
\item Stage 9 packaging/implementation, which consumes Stage 3 as the scheduling doctrine for deployable execution.
\end{itemize}

\paragraph*{5.3 Crosswalk hooks}

Stage 3 uses and \textbf{MUST} keep stable the following join keys and scheduling references:

\begin{itemize}
\item Phase IDs and phase week markers
\item Protocol cards \textbf{C1}–\textbf{C18}
\item Metric IDs: \textbf{MET-2B}-###
\item Domain tags: \textbf{RUN}, \textbf{SWIM}, \textbf{CAL}, \textbf{RUCK}, \textbf{STR}, \textbf{BODYCOMP}, \textbf{DURABILITY}, \textbf{RECOVERY}, \textbf{AER}
\item Environment Unit \textbf{ID} concept and tool consistency concept (Stage 2.E)
\end{itemize}

\subsubsection*{6) Gate / Acceptance Criteria}

Stage 3 is complete and deployable when:

\begin{itemize}
\item Every phase has protected mid-phase and end-phase test windows consistent with Stage 3 doctrine and compatible with Stage 1.F admissibility posture and Stage 2 replication rules.
\item Macro stress doctrine (work-to-deload rhythm) is explicit and implementable without prescribing workouts.
\item Anti-stacking/spacing doctrine is explicit and prevents:
  \begin{itemize}
  \item fatigue contamination,
  \item unsafe escalation,
  \item and cross-domain interference (ROCK is not used; \textbf{RUCK} and \textbf{RUN} remain distinct).
  \end{itemize}
\item Environment reslot/substitution posture is explicit for Grand Forks constraints and preserves safety and validity:
  \begin{itemize}
  \item indoor substitution only when comparable under Stage 2 validity,
  \item otherwise reslot.
  \end{itemize}
\item Reslot doctrine is explicit for:
  \begin{itemize}
  \item missed tests,
  \item invalid tests,
  \item illness,
  \item regression,
  \item and risk flags.
  \end{itemize}
\item Replication enablement posture is explicit and achievable without overtesting.
\item No contradictions exist with Stage 0.5 water governance and Stage 1.F stop posture:
  \begin{itemize}
  \item underwater/breath-hold is never scheduled without governance prerequisites,
  \item any breach forces conservative posture.
  \end{itemize}
\item Stage 3 deliverables 3.A–3.E are scoped and interlocked without introducing new tests, thresholds, domains, or identifiers.
\end{itemize}

\subsubsection*{7) Common Failure Modes + Controls}

\begin{itemize}
\item Overtesting and fatigue contamination → enforce deload/test window protection; cap high-cost test density; require spacing.
\item Missed tests “made up” by stacking → reslot discipline; no stacking; no compression after missed weeks.
\item Unsafe environment testing (ice/heat/air quality/low visibility) → postpone or substitute indoors only when Stage 2 validity/logging preserves comparability; otherwise reslot.
\item Novelty near gates (new route/tool/shoes/caffeine/timekeeping) → freeze posture for verification windows; log any forced changes as comparability breaks; defer gate-grade testing until stabilized.
\item Insufficient replication opportunities → recurring protected windows structured to enable replication across windows, not within a single week.
\item Tool or route drift → enforce Environment Unit IDs and tool-change logging per Stage 2.E; treat drift as a comparability break until stabilized and replicated.
\item Water governance breaches → supervised-only posture; underwater is not scheduled without governance prerequisites; any breach is inadmissible.
\item Pool access volatility causing invalid \textbf{SWIM} evidence → pool verification mandatory; reslot \textbf{SWIM} tests rather than substituting non-swim evidence.
\item Long-duration tests adjacent to other major tests → spacing doctrine; separate high-cost events by minimum spacing intervals.
\item Invalid tests used anyway to support gates → \textbf{INVALID} provides no gate credit; require retest in next protected window; gate posture defaults \textbf{HOLD} for that metric until valid evidence exists.
\item Compression after illness/injury → reslot and conservative return posture; do not force tests under unresolved risk flags.
\item Gate claims made without replication → replication tracking required; default \textbf{HOLD} (Replication) for that metric until replication exists.
\end{itemize}

\subsubsection*{8) Versioning Notes}

\begin{itemize}
\item Frozen scheduling doctrine primitives:
  \begin{itemize}
  \item protected deload/test windows posture,
  \item macro stress rhythm posture (work-to-deload),
  \item anti-stacking/spacing doctrine,
  \item reslot posture for invalid/missed tests,
  \item environment-driven substitution limits for Grand Forks conditions,
  \item replication enablement posture (cadence sufficient to achieve 2 of last 3 \textbf{VALID} without overtesting).
  \end{itemize}
\item Refinable only via logged change:
  \begin{itemize}
  \item phase-specific window placement refinements,
  \item environment substitution clarifications that do not alter Stage 2 validity rules,
  \item minor calendar structure clarifications that do not alter governance meaning.
  \end{itemize}
\item Patch/change control posture:
  \begin{itemize}
  \item any change to Stage 3 scheduling doctrine \textbf{MUST} be logged,
  \item reviewed for downstream impact (Stages 4–5),
  \item and checked for Stage 8 crosswalk implications before treated as deployable.
  \end{itemize}
\end{itemize}

\subsection*{Stage 3 Contents Map}
\phantomsection
\addcontentsline{toc}{subsection}{Stage 3 Contents Map}

\begin{itemize}
\item 3.A Design Principles for Testing and Deloads — Defines scheduling principles that protect validity, safety, and comparability for all tests across Phase 00–Phase 13, including protected window and anti-stacking doctrine.
\item 3.B Phase-Level Stress and Test Map — Establishes phase-level placement of protected test windows and stress rhythm without prescribing training, including mid-phase and end-phase window logic and minimum spacing rules.
\item 3.C Test Cadence by Metric Type — Defines eligibility cadence for each metric category (Gate-Critical performance tests vs supporting markers vs trends) to enable replication without chronic testing burden.
\item 3.D Integrated Phase Calendars Compact Form — Provides compact phase-level calendars showing protected test windows and gate points without day-by-day programming, using Phase Week markers and window labels only.
\item 3.E Adjustment Rules for Illness, Regression, and Risk Flags — Defines reslot and regression scheduling rules when safety, validity, or feasibility constraints disrupt planned windows, including invalid-test retest posture and environment-driven reslots.
\end{itemize}

\end{document}